\section{Linjär algebra}

\subsection{Vektorer}
\subsubsection{Skrivsätt}
Utöver koordinat-betäckningen $\vec{v}=(v_1,\dots,v_n)$, kan vektorer även beskrivas
som matriser (se \ref{s:matris}) i antingen \emph{kolonnvektor} form som en $n\times 1$ matris:
$$
	\vec{v}=\qty(\begin{matrix}
		v_1\\\vdots\\v_n
	\end{matrix})
$$
eller \emph{radvektor} form som en $1\times n$ matris:
$$
	\vec{v}=\underbrace{\qty(v_1,\dots,v_n)}_{\text{\makebox[0pt]{Ofta $[v_1,\dots,v_n]$ för att undvika förväxling med koordinater}}}
$$

\subsubsection{Vektoroperationer}

\paragraph{\href{https://sv.wikipedia.org/wiki/Skalärprodukt}{Skalärprodukter}:}

\begin{features}
	\textbf{Egenskaper av skalärprodukten}

	Skalärprodukten är bilinjär, detta innebär:
	\begin{itemize}
		\item \textbf{Skalbar:} $a\vec{u}\cdot\vec{v}=\vec{u}\cdot a\vec{v}=a(\vec{u}\cdot\vec{v})$
		\item \textbf{Distributivitet:} $\vec{u}\cdot(\vec{v}+\vec{w})=\vec{u}\cdot\vec{v}+\vec{u}\cdot\vec{w}$
		\item \textbf{Kommutativitet:} $\vec{u}\cdot\vec{v}=\vec{v}\cdot\vec{u}$
	\end{itemize}
\end{features}

\paragraph{\href{https://sv.wikipedia.org/wiki/Kryssprodukt}{Kryssprodukten}:}

\begin{features}
	\textbf{Egenskaper av kryssprodukten}

	\begin{itemize}
		\item \textbf{Skalbar:} $a\vec{u}\cross\vec{v}=\vec{u}\cross a\vec{v}=a(\vec{u}\cross\vec{v})$
		\item \textbf{Distributivitet:} $\vec{u}\cross(\vec{v}+\vec{w})=\vec{u}\cross\vec{v}+\vec{u}\cross\vec{w}$
		\item \textbf{Anti-kommutativitet:} $\vec{u}\cross\vec{v}=-\vec{v}\cross\vec{u}$
	\end{itemize}
\end{features}

\subsubsection{Underrum}

$\spn(\vec{v}_1,\dots,\vec{v}_n)\subsetneq\mathbb{R}^n$ kallas för ett underrum
till $\mathbb{R}^n$. Altså är $\vec{v}_1,\dots,\vec{v}_n$ vektorer som inte spänner
upp $\mathbb{R}^n$.

\subsubsection{Dimension}
Dimensionen av ett underrum $W$, skriven $\dim W$, är antalet vektorer i en bas
för underrummet, vilket motsvarar det maximala antalet linjärt oberoende vektorer
i rummet.

\begin{ber}
	\textbf{Exempel på icke-trivell dimension beräkning:}

	Låt $\Pi:x+2y-z=0$. För att beräkna $\dim\Pi$, måste antalet linjärt
	oberoende vektorer krävda för att spänna upp $\Pi$ beräknas. Vi skriver om i
	vektorform, genom att låta $y=s,z=t\in\mathbb{R}$:
	\begin{align*}
		&x+2y-z=0\implies x=t-2s\\
		&\therefore\qty(\begin{matrix}
			x\\y\\z
		\end{matrix})=\qty(\begin{matrix}
			t-2s\\s\\t
		\end{matrix})=s\qty(\begin{matrix}
			-2\\1\\0
		\end{matrix})+t\qty(\begin{matrix}
			1\\0\\1
		\end{matrix})\\
		&\qty(\begin{matrix}
			-2&1&0\\
			1&0&1
		\end{matrix})\sim\qty(\begin{matrix}
			\boxed{-2}&1&0\\
			0&\boxed{1}&2
		\end{matrix}),\quad\text{pivotelement i varje rad}\implies\text{linjärt oberoende}\\
		&\therefore\spn\qty{\qty(\begin{matrix}
			-2\\1\\0
		\end{matrix}),\qty(\begin{matrix}
			1\\0\\1
		\end{matrix})}=\Pi\implies\dim\Pi=2
	\end{align*}
\end{ber}

\subsection{Linjer}

\subsubsection{Vektorform}\label{s:linalg:linje:vektor-form}
En linje i planet $\mathbb{R}^2$ eller rummet $\mathbb{R}^3$ i vektorform är en linje $L$ skriven i formen:
$$
	L:P+\vec{v}t,\quad t\in\mathbb{R}
$$
Där $P\in L$ är en punkt på linjen, och $\vec{v}$ riktningsvektorn av linjen.

\begin{ber}
\textbf{Beräkning från två angivna punkter $P$ och $Q$}

Vi vet nu triviellt en punkt på linjen ($P$ eller $Q$).
Riktningsvektorn är sedan lätt beräknad genom att finna vektorn mellan punkterna:
$$
	\vec{v}=\overrightarrow{PQ}\quad(\text{eller }\overrightarrow{QP})
$$
Ekvationen för linjen blir sedan:
$$
	L:P+\vec{v}t,\quad t\in\mathbb{R}
$$
\end{ber}

\begin{omv}
\textbf{Omvandling till affin form (\ref{s:linalg:linje:affin-form})}

Given linjen i planet $\mathbb{R}^2$ (processen för rummet är analogt):
$$
	L:(A,B)+t(C,D),\quad t\in\mathbb{R}
$$
\end{omv}

\subsubsection{Affin form}\label{s:linalg:linje:affin-form}
En linje i planet $\mathbb{R}^2$ (OBS, inte rummet!) i affin form är en linje $L$ skriven i formen:
$$
	L=ax+by=c
$$
där $a,b,c\in\mathbb{R}$ är konstanter.


\begin{der}
	\textbf{Derivering av den affina formen av linjer}

	Se \textbf{Derivering av den affina formen av plan} i \ref{s:linalg:plan:affin-form},
	processen är analog för linjer i planet.
\end{der}

\begin{omv}
\textbf{Omvandling till vektorform (\ref{s:linalg:linje:vektor-form})}

Given en linje av form $L:ax+by=c$, låt en av variablerna vara $t$.
Till exempel: $x\leteq t^*$. Lös sedan för den andra variabeln:
\begin{align*}
	at+by=c\implies by&=c-at^*\\
	\therefore y&=\frac{c-at^*}{b}=\frac{c}{b}-\frac{a}{b}t^*
\end{align*}
Alternativt, genom att sedan definera $t=bt^*$ får vi ut en lite enklare formel:
$$
	y=\frac{c}{b}-at
$$
Vektorformsekvationen blir sedan:
$$
	L:\qty(0,\frac{c}{b})+t(b,-a),\quad t\in\mathbb{R}
$$
För ett finare uttryck än $\qty(0,\frac{c}{b})$, om du vet någon punkt $\qty(\alpha,\beta)$
ligger på linjen, så kan den användas istället.
\end{omv}

\subsection{\href{https://sv.wikipedia.org/wiki/Plan_(geometri)}{Plan}}

\subsubsection{Vektorform}

En ett plan i rummet $\mathbb{R}^3$ i vektorform är ett plan $\Pi$ skriven i formen:
$$
	\Pi:P+\vec{u}s+\vec{v}t,\quad s,t\in\mathbb{R}
$$
Där $P\in\Pi$ är en punkt i planet, och $\vec{u},\vec{v}$ riktningsvektornerna av linjen.

\subsubsection{Affin form}\label{s:linalg:plan:affin-form}

Ett plan i affin form är ett plan $\Pi$ skriven i formen:
$$
	\Pi:ax+by+cz=d
$$
där $a,b,c,d\in\mathbb{R}$ är konstanter.

\begin{der}
\textbf{Derivering av den affina formen av plan}

Låt $O$ vara en punkt på planet.
Vektorn mellan $O$ varje annan punkt på planet, $P$, måste vara vinkelrät mot normalvektorn,
eftersom normalvektorn är vinkelrät mot planet, och $\overrightarrow{OP}$ är en vektor i planet.

Därför kan vi skriva:
\begin{align*}
	\vec{n}\cdot\overrightarrow{OP}&=0\\
	\implies \vec{n}\cdot(P-O)&=0\\
	\therefore \vec{n}\cdot P&=\vec{n}\cdot O
\end{align*}

Eftersom vi redan, per definition av frågan, känner till värdet av $\vec{n}\leteq(a,b,c)$
och $O$, så blir $\vec{n}\cdot O$ någon konstant, som vi betäcknar $d$.
Låt slutligen våran punkt på planet $P\leteq(x,y,z)$:
\begin{align*}
	(a,b,c)\cdot(x,y,z)&=d\\
	\therefore ax+by+cz&=d
\end{align*}
\end{der}

\begin{ber}
	Given en punkt på planet, och en normal vektor, se \textbf{Derivering av den affina formen av plan}.
\end{ber}

\subsection{Avstånd}

\subsection{\href{https://sv.wikipedia.org/wiki/Bas_(linjär_algebra)}{Baser}}

\subsubsection{Bildning av baser}

\paragraph{Matriskontroll:}
För att kontrollera om vektorerna $\{\vec{v}_j\}_{j=1}^p$ bildar en bas,
bilda en matris med vektorerna som kolumner och betrakta om dess determinant är nollskiljt:
$$
	\qty|\begin{matrix}
		\vec{v}_1^\top,\vec{v}_2^\top,\dots,\vec{v}_p^\top
	\end{matrix}|\stackrel{?}{=}0
$$
Om den är det, så bildar vektorerna en bas, annars gör de inte det.
Detta följer då en determinant av $0$ innebär att vektorerna är linjärt beroende.

\subsection{Matriser}\label{s:matris}
En matris, ofta betecknade $\matr{A}$ och $\matr{B}$, är ett 2 dimensionellt
rutnät/tabell av tal $a_{ij}$, där $i$ och $j$ är dess rad respektive kolumn.

Matriser är utskrivna som:
$$
\matr{A}=\qty(\begin{matrix}
	a_{11}&a_{12}&\dots&a_{1n}\\
	a_{21}&a_{22}&\dots&a_{2n}\\
	\vdots&\vdots&\ddots&\vdots\\
	a_{m1}&a_{m2}&\dots&a_{mn}
\end{matrix})\in\mathbb{R}^{m\times n},
$$
där $\matr{A}$ sägs vara av \emph{typ} eller \emph{storlek} $m\times n$. 
Alternativt kan $\matr{A}$ skrivas i kortform som $\matr{A}=(a_{ij})$.

Matriser dyker upp ofta naturligt i flera sammanhang, bland annat då vi löser
LES med gausselimination. Dessa skrivs då som \emph{augmenterade matriser}:
$$
	\left\{\,\begin{matrix}
		ax&+&by&=&c\\
		dx&+&ey&=&f
	\end{matrix}\right.\quad\longrightarrow\quad\begin{amatrix}{2}
		a&b&c\\
		d&e&f
	\end{amatrix}
$$
där det räta sträcket representerar ett $=$ tecken.

\subsubsection{Matrisoperationer}

\paragraph{Addition:}
För matriserna $\matr{A}=(a_{ij})$ och $\matr{B}=(b_{ij})$, båda av samma typ,
defineras $\matr{A}+\matr{B}$ som:
$$
	\matr{A}+\matr{B}\defeq(a_{ij}+b_{ij})=\qty(\begin{matrix}
		a_{11}+b_{11}&a_{12}+b_{12}&\dots&a_{1n}+b_{1n}\\
		a_{21}+b_{21}&a_{22}+b_{22}&\dots&a_{2n}+b_{2n}\\
		\vdots&\vdots&\ddots&\vdots\\
		a_{m1}+b_{m1}&a_{m2}+b_{m2}&\dots&a_{mn}+b_{mn}
	\end{matrix})
$$

\begin{features}
	\textbf{Egenskaper av matrisaddition}

	På grund av dess definition, har matrisaddition samma egenskaper som skaläraddition.
\end{features}

\paragraph{Multiplikation med skalärer:}
Om $\matr{A}\in\mathbb{R}^{m\times n}$ och
$\lambda\in\mathbb{R}$ så är dess produkt:
$$
	\lambda\matr{A}\defeq(\lambda a_{ij})=\qty(\begin{matrix}
		\lambda a_{11}&\lambda a_{12}&\dots&\lambda a_{1n}\\
		\lambda a_{21}&\lambda a_{22}&\dots&\lambda a_{2n}\\
		\vdots&\vdots&\ddots&\vdots\\
		\lambda a_{m1}&\lambda a_{m2}&\dots&\lambda a_{mn}
	\end{matrix})
$$

\begin{features}
	\textbf{Egenskaper av multiplikation av matriser med skalärer}

	På grund av dess definition, har multiplikation av matriser med skalärer samma
	egenskaper som skalärmultiplikation.
\end{features}

\paragraph{Multiplikation:} Om $\matr{A}$ är av typ $m\times n$, och $B$ av $n\times p$
(\textbf{OBS} andra respektive första dimensionen, i detta fall $n$, av $\matr{A},\matr{B}$ måste vara lika!),
så är $\matr{AB}$ definerat som (för definition av $(\cdot)^\top$, se \ref{s:matris-transponering}):
\begin{align*}
	\matr{AB}&\defeq\qty(\begin{matrix}
		\vec{\alpha}_1&
		\dots&
		\vec{\alpha}_i&
		\dots&
		\vec{\alpha}_m
	\end{matrix})\qty(\begin{matrix}
		\vec{\beta}_1^\top\\
		\vdots\\
		\vec{\beta}_j^\top\\
		\vdots\\
		\vec{\beta}_n^\top
	\end{matrix})\\
	&=\qty(\begin{matrix}
		a_{11}&\dots&a_{1n}\\
		\vdots&&\vdots\\
		a_{i1}&\dots&a_{in}\\
		\vdots&&\vdots\\
		a_{m1}&\dots&a_{mn}
	\end{matrix})\qty(\begin{matrix}
		b_{11}&\dots&b_{1j}&\dots&b_{1p}\\
		\vdots&&\vdots&&\vdots\\
		b_{n1}&\dots&b_{nj}&\dots&b_{np}\\
	\end{matrix})\\
	&=\qty(\begin{matrix}
		\vec{\alpha}_1\cdot\vec{\beta}_1&\cdots&\vec{\alpha}_1\cdot\vec{\beta}_j&\cdots&\vec{\alpha}_1\cdot\vec{\beta}_n\\
		\vdots&&\vdots&&\vdots\\
		\vec{\alpha}_i\cdot\vec{\beta}_1&\cdots&\vec{\alpha}_i\cdot\vec{\beta}_j&\cdots&\vec{\alpha}_i\cdot\vec{\beta}_n\\
		\vdots&&\vdots&&\vdots\\
		\vec{\alpha}_m\cdot\vec{\beta}_1&\cdots&\vec{\alpha}_m\cdot\vec{\beta}_j&\cdots&\vec{\alpha}_m\cdot\vec{\beta}_n
	\end{matrix}),
\end{align*}
eller, alternativt uttryckt som:
$$
	(\matr{AB})_{ij}=\sum_{k=1}^na_{ik}b_{kj}.
$$

\begin{features}
	\textbf{Egenskaper av matrismultiplikation}

	\textbf{OBS} matrismultiplikation är \emph{ej} kommutativ!

	\begin{itemize}
		\item \textbf{Associativitet:} $\matr{A}(\matr{B}\matr{C})=(\matr{A}\matr{B})\matr{C}$ \begin{proof}
			Låt $\matr{A},\matr{B},\matr{C}$ vara av typer $m\times n$, $n\times p$, $p\times q$.
			\begin{align*}
				(\overbrace{(\matr{A}\matr{B})}^{m\times p}\matr{C})_{ij}
				&\defeqto\sum_{k=1}^p(\matr{A}\matr{B})_{ik}\matr{C}_{kj}
				=\sum_{k=1}^p\qty(\sum_{\ell=1}^na_{i\ell}b_{\ell k})c_{kj}\\
				&=\sum_{k=1}^p\sum_{\ell=1}^na_{i\ell}b_{\ell k}c_{kj}
				=\sum_{\ell=1}^n\qty(\sum_{k=1}^pa_{i\ell}b_{\ell k}c_{kj})\\
				&=\sum_{\ell=1}^na_{i\ell}\qty(\sum_{k=1}^pb_{\ell k}c_{kj})
				=\sum_{\ell=1}^na_{i\ell}(\matr{B}\matr{C})_{\ell j}
				=(\matr{A}(\matr{B}\matr{C}))_{ij}
			\end{align*}
		\end{proof}
		\item \textbf{Distributivitet:} $\matr{A}(\matr{B}+\matr{C})=\matr{A}\matr{B}+\matr{A}\matr{C}$ 
	\end{itemize}
\end{features}

\paragraph{Transponering:}\label{s:matris-transponering}
transponatet av en matris $\matr{A}$, skriven $\matr{A}^\top$,
byter plats på matrisens rader och kolumner (\textbf{OBS} inte en $90^\circ$ rotation).

\begin{ber}
	\textbf{Exempel på transponering:}

	Låt $\matr{A}$ vara sådan att:
	$$
		\matr{A}=\qty(\begin{matrix}
			\vec{a}_1^\top\\
			\vdots\\
			\vec{a}_m^\top
		\end{matrix})\in\mathbb{R}^{m\times n}
	$$, där $\vec{a}^\top_j\in\mathbb{R}^n, j\in[1,m]\cap\mathbb{Z}$
	är radvektorer. Då gäller att dess transponat beräcknas enligt:
	$$
		\matr{A}^\top=\qty(\begin{matrix}
			\vec{a}_1&\dots&\vec{a}_m
		\end{matrix})\in\mathbb{R}^{n\times m}
	$$
\end{ber}

\begin{features}
	\textbf{Egenskaper av transponatet}

	\begin{itemize}
		\item \textbf{Skalbar:} $(\lambda\matr{A})^\top=\lambda\matr{A}^\top$
		\item \textbf{Involution:} $\qty(\matr{A}^\top)^\top=\matr{A}$
		\item \textbf{Additiv:} $(\matr{A}+\matr{B})^\top=\matr{A}^\top+\matr{B}^\top$
		\item \textbf{Anti-homomorfism:} $(\matr{A}\matr{B})^\top=\matr{B}^\top\matr{A}^\top$
	\end{itemize}
\end{features}

\subsubsection{Särskilda matriser}

\paragraph{Nollmatrisen:} Nollmatrisen skrivs $\matr{0}$ och är definerad som en
$m\times n$ matris där $\matr{0}_{ij}=0\forall i,j$.

För Nollmatrisen och någon annan matris $\matr{A}$ av godtycklig typ och skalär $\lambda$
gäller följande räckneregler:
\begin{itemize}
	\item $\matr{0}\matr{A}=\matr{A}\matr{0}=\matr{0}$
	\item $\matr{A}+\matr{0}=\matr{A}$
	\item $\matr{A}+(-1)\matr{A}=\matr{0}$
	\item $\lambda\matr{0}=\matr{0}$
\end{itemize}

\paragraph{Enhetsmatrisen:} Enhetsmatrisen, också känd som identitetsmatrisen,
betecknad $\matr{E}_n$ (eller bara $\matr{E}$ om kontext gör storleken meningslös),
är en matris definerad enligt:
\begin{align*}
	&\matr{E}_1=\qty(\begin{matrix}
		1
	\end{matrix}),&\matr{E}_2=\qty(\begin{matrix}
		1 & 0\\
		0 & 1
	\end{matrix})&\\
	&\matr{E}_3=\qty(\begin{matrix}
		1 & 0 & 0\\
		0 & 1 & 0\\
		0 & 0 & 1
	\end{matrix}),&\matr{E}_n=\underbrace{\qty(\begin{matrix}
		1 & 0 & 0 & \dots & 0\\
		0 & 1 & 0 & \dots & 0\\
		0 & 0 & 1 & \dots & 0\\
		\vdots & \vdots & \vdots & \ddots & \vdots\\
		0 & 0 & 0 & 0 & 1
	\end{matrix})}_\text{$n$ ggr}&
\end{align*}
d.v.s. en matris av typ $n\times n$, som är noll vid alla positioner utom dess
diagonal.

Dess namn följer av dess egenskaper som \emph{enheten} av en $n\times n$ matrisring.
Därmed följer även att för någon annan matris $\matr{A}$ så är $\matr{E}\matr{A}=\matr{A}\matr{E}=\matr{A}$.

\paragraph{Symetriska matriser:} Om det gäller för en matris $\matr{A}$ att $\matr{A}^\top=\matr{A}$ (för definition
av $(\cdot)^\top$, se \ref{s:matris-transponering}), så kallas $\matr{A}$ symmetrisk.

\subsubsection{Kärnor}
Mängden av alla lösningar till matrisekvationen $\matr{A}\vec{x}=\vec{0}$ kallas
för $\matr{A}$:s kärna (eller nollrum) och betecknas $\ker(\matr{A})$ eller $\noll(\matr{A})$.
Denna defineras enligt:
$$
	\ker(\matr{A})=\qty{\vec{x}\in\mathbb{R}^n:\matr{A}\vec{x}=\vec{0}}
$$
där $\matr{A}$ är av typ $m\times n$.

\begin{ber}
	\textbf{Kärnan av vissa särskilda matriser:}

	\begin{itemize}
		\item $\ker(\matr{0}_{m,n})=\mathbb{R}^n$
		\item $\ker(\matr{E}_n)=\qty{\vec{0}}$
	\end{itemize}
\end{ber}

\subsubsection{Nolldimensioner}
Nolldimensionen av en matris $\matr{A}$ är dimensionen av dess kärna,
$$
	\underbrace{\nullity(\matr{A})}_\text{\text{\makebox[0pt]{På svenska mer känd som $\nolldim(\matr{A})$}}}\defeq\dim(\ker\matr{A})
$$

\subsubsection{Kolonnrum}\label{s:matris:kolonnrum}
Kolonnrummet för en matris $\matr{A}\in\mathbb{R}^{m\times n}$, betecknat $\koll(\matr{A})$,
är det underrum i $\mathbb{R}^m$ som spänns upp av matrisens kolonnvektorer. Dimensionen
för kolonnrummet, $\dim(\koll\matr{A})$, kallas för matrisens \emph{rang} och betecknas
$\rang(\matr{A})$. Vid gausselimination motsvarar rangen antalet pivotelement
i trappstegsformen.

\subsubsection{Dimensionssatsen}\label{s:matris:dimensionssatsen}
Dimensionssatsen, ofta kallad rangsatsen, relaterar dimensionerna för en matris
kolonnrum och nollrum med dess totala antal kolonner.

Om $\matr{A}\in\mathbb{R}^{m\times n}$ (altså med $n$ kolonner) så gäller:
$$
	\dim(\koll\matr{A}) + \nullity(\matr{A})=n
$$

\begin{features}
	\textbf{Betydelse i trappstegsform}
	
	Satsen ger för linjära ekvationssystem att varje kolonn i en radreducerad
	matris måste antingen vara en pivotkolonn eller motsvara en fri variabel.
	\begin{itemize}
		\item $\dim(\koll\matr{A}) = \text{Antal pivotkolonner (bundna variabler)}$
		\item $\dim(\nullity\matr{A}) = \text{Antal fria variabler}$
	\end{itemize}
	Eftersom en kolonn inte kan vara båda samtidigt, blir summan alltid lika med
	det totala antalet kolonner $n$.
\end{features}

\begin{ber}
	\textbf{Exempel på beräkning enligt dimensionssatsen:}

	Betrakta $\matr{A}\in\mathbb{R}^{3\times 4}$:
	$$
		\matr{A}=\qty(\begin{matrix} 
			1&-2&3&1\\ 
			2&-4&7&4\\ 
			3&-6&10&5 
		\end{matrix})
	$$
	Efter gausselimination får vi ut, i reducerad trappstegsform, matrisen $\matr{R}$:
	$$
		\matr{A}\sim\matr{R}=\qty(\begin{matrix} 
			\boxed{1} & -2 & 0 & -5 \\ 
			0 & 0 & \boxed{1} & 2 \\ 
			0 & 0 & 0 & 0 
		\end{matrix})
	$$
	
	\textbf{Tolkning av trappstegsformen $\matr{B}$:}
	\begin{itemize}
		\item Kolonn 1 och 3 innehåller pivotelement. Detta betyder att de
		ursprungliga kolonnerna 1 och 3 i $\matr{A}$ bildar en bas för kolonnrummet.
		Rangen är alltså:
		$$
			\rang(\matr{A})=\dim(\koll\matr{A})=2
		$$
		
		\item Kolonn 2 och 4 saknar pivotelement. Detta innebär att variablerna
		$x_2$ och $x_4$ är fria variabler när vi löser det homogena systemet $\matr{A}\vec{x}=\vec{0}$.
		Eftersom varje fri variabel ger en unik basvektor i $\ker(\matr{A})$:
		$$
			\nullity(\matr{A})=\dim(\ker\matr{A})=2
		$$
	\end{itemize}

	Slutligen, enligt dimensionssatsen bör summan av $\rang(\matr{A})$ och
	$\nullity(\matr{A})$ vara antalet kolloner, $4$:
	$$
		\rang(\matr{A})+\nullity(\matr{A})=2+2=4\quad\checkmark
	$$
\end{ber}
